\chapter[2016 --- Ohsumi]{2016: Yoshinori Ohsumi}
\label{ch:2016_yoshinoriohsumi}

\section*{Main Contribution}

\textit{Discovered the mechanisms of autophagy---how cells recycle their own components---a fundamental process essential for health and implicated in cancer, neurodegeneration, and aging.}

\section{Official Citation}

for his discoveries of mechanisms for autophagy

\section{Background and Historical Context}

Yoshinori Ohsumi, a Japanese cell biologist born in 1945 in Fukuoka, Japan, was awarded the 2016 Nobel Prize in Physiology or Medicine ``for his discoveries of mechanisms for autophagy.'' Autophagy, derived from the Greek words for ``self'' and ``eating,'' is a fundamental cellular process by which cells degrade and recycle their own components. Ohsumi's work, conducted at the Tokyo Institute of Technology, elucidated the molecular machinery of autophagy and revealed its essential roles in cellular physiology, disease, and survival.

The concept of autophagy was first described in the 1960s, when Christian de Duve and colleagues observed the presence of double-membrane-bound organelles containing cellular debris in rat liver cells. These organelles, which were named ``autophagosomes,'' appeared to sequester cytoplasmic components and deliver them to lysosomes for degradation. Subsequent electron microscopy studies confirmed the existence of autophagy in a wide range of cell types and organisms, and suggested that the process played important roles in cellular remodeling, starvation responses, and tissue homeostasis.

However, the molecular mechanisms underlying autophagy---the genes and proteins responsible for the formation of autophagosomes, their fusion with lysosomes, and the regulation of the process in response to various cellular signals---remained largely unknown for decades. The identification and characterization of these mechanisms became one of a major challenges in cell biology, with significant implications for our understanding of health and disease.

Autophagy was known to be induced under conditions of nutrient deprivation, suggesting that it served as a survival mechanism by providing cells with essential nutrients and energy during periods of starvation. However, the full scope of autophagy's functions was not appreciated until the molecular machinery was elucidated, revealing its involvement in virtually every aspect of cellular physiology, including development, differentiation, immune defense, and the prevention of neurodegenerative diseases and cancer.

\section{The Discovery/Contribution}

Ohsumi's contributions to the understanding of autophagy began in the early 1990s, when he was working at the National Institute for Basic Biology in Okazaki, Japan. He chose to study autophagy in baker's yeast (\textit{Saccharomyces cerevisiae}) because of its genetic tractability and the availability of powerful genetic tools. While autophagy was well known to occur in mammalian cells, its occurrence in yeast had not been demonstrated.

Ohsumi developed a genetic screen to identify yeast mutants that were defective in autophagy. His approach was based on a simple but elegant logic: if autophagy was essential for the recycling of cellular components during starvation, then mutants lacking functional autophagy would be unable to survive prolonged nutrient deprivation. By subjecting yeast cells to nitrogen starvation and screening for colonies that failed to survive, Ohsumi identified a collection of autophagy-defective (``\textit{atg}'') mutants.

In a landmark study published in 1993, Ohsumi identified the first 15 genes essential for autophagy in yeast, naming them \textit{APG1} through \textit{APG15} (later renamed \textit{ATG} genes for ``autophagy-related''). He demonstrated that these genes were required for the formation of autophagosomes and that their inactivation led to a block in the autophagy pathway. The identification of these genes provided the first molecular handles on the autophagy machinery and opened up the field of autophagy research.

Subsequent work by Ohsumi and many other investigators elucidated the function of the ATG gene products and revealed the molecular mechanisms of autophagy in remarkable detail. Ohsumi's laboratory identified the ATG proteins and characterized their functions, establishing the core machinery of the autophagy pathway. The ATG genes were found to encode a diverse array of proteins, including kinases, ubiquitin-like conjugation systems, lipid-modifying enzymes, and membrane-tethering factors. Together, these proteins orchestrated the formation of the autophagosome, a double-membrane-bound organelle that sequestered portions of the cytoplasm and delivered them to the lysosome for degradation.

The autophagy pathway was shown to involve several distinct steps:

\textbf{Initiation:} Autophagy is initiated in response to nutrient deprivation or other cellular stresses. The kinase TOR (target of rapamycin), a central regulator of cell growth and metabolism, is inhibited under conditions of nutrient limitation, leading to the activation of the autophagy-initiating kinase complex consisting of ATG1, ATG13, and ATG17. Under nutrient-rich conditions, TOR phosphorylates ATG13 and inhibits its activity, preventing autophagy. When nutrients are scarce, TOR is inactivated, ATG13 is dephosphorylated and activated, and the ATG1 kinase complex initiates the autophagy cascade. This TOR-dependent regulation ensures that autophagy is active only when nutrients are limiting and the cell needs to recycle its own components for energy and building blocks.

\textbf{Nucleation:} The activated ATG1 kinase complex recruits a second complex consisting of ATG6, ATG14, Vps34 (a class III phosphatidylinositol 3-kinase), and Vps15 to the site of autophagosome formation, known as the phagophore assembly site (PAS). Vps34 generates phosphatidylinositol 3-phosphate (PI3P), a lipid second messenger that recruits additional ATG proteins to the PAS and promotes the expansion of the phagophore membrane. The PI3P produced by Vps34 is recognized by WIPI proteins and DFCP1, which serve as PI3P effectors that recruit downstream components of the autophagy machinery to the phagophore.

\textbf{Elongation:} The phagophore membrane expands by the sequential addition of lipid and protein components. Two ubiquitin-like conjugation systems play essential roles in this process. In the first system, ATG12 is conjugated to ATG7 and then to ATG10, forming an ATG12-ATG5-ATG16 complex that associates with the phagophore membrane. In the second system, the cytosolic protein ATG8 (LC3 in mammals) is cleaved by the protease ATG4 to expose a C-terminal glycine residue, which is then conjugated to the lipid phosphatidylethanolamine (PE) by the sequential action of ATG7 and ATG3. ATG8-PE decorates the phagophore membrane and plays essential roles in membrane expansion, cargo recognition, and autophagosome closure.

\textbf{Closure:} The phagophore membrane eventually closes around the cytoplasmic cargo to form a complete double-membrane vesicle---the autophagosome. The mechanisms controlling autophagosome closure are still being elucidated but involve the ESCRT (endosomal sorting complexes required for transport) machinery, which mediates membrane scission events in other cellular processes.

\textbf{Fusion:} The autophagosome fuses with a lysosome to form an autolysosome. This fusion process requires SNARE proteins, Rab GTPases, and other components of the vesicle trafficking machinery. The Rab7 GTPase plays a particularly important role in promoting autophagosome-lysosome fusion, while the SNARE proteins VAMP8, STX17, and SNAP29 mediate the membrane fusion event. The acidic environment (pH approximately 4.5--5.0) and hydrolytic enzymes of the lysosome degrade the contents of the autophagosome.

\textbf{Recycling:} The degradation products are exported from the autolysosome to the cytoplasm by permeases and other transporters, where they can be used for biosynthesis and energy production. The amino acids, fatty acids, nucleotides, and other small molecules released by lysosomal degradation provide the raw materials for new cellular components and can be oxidized by mitochondria to generate ATP during periods of nutrient deprivation.

Ohsumi's genetic approach, combined with the biochemical and cell biological studies that followed, provided a comprehensive picture of the autophagy pathway and revealed its remarkable complexity and versatility. The identification of ATG genes in yeast, and the subsequent discovery of their conserved homologs in mammals, demonstrated that the autophagy machinery was highly conserved across eukaryotes, underscoring its fundamental importance for cellular life.

\section{Impact on Modern Medicine}

The impact of Ohsumi's discoveries on modern medicine has been significant. The elucidation of the molecular machinery of autophagy has revolutionized our understanding of cellular physiology and has provided critical insights into the pathogenesis of a wide range of human diseases.

\textbf{Neurodegenerative diseases:} Autophagy plays a crucial role in the removal of misfolded and aggregated proteins from neurons, and defects in autophagy have been implicated in the pathogenesis of Alzheimer's disease, Parkinson's disease, Huntington's disease, and amyotrophic lateral sclerosis (ALS). In Alzheimer's disease, impaired autophagy leads to the accumulation of amyloid-beta plaques and tau tangles, the characteristic pathological hallmarks of the disease. In Parkinson's disease, defects in autophagy contribute to the accumulation of alpha-synuclein aggregates and the loss of dopaminergic neurons. The activation of autophagy has been explored as a therapeutic strategy for these and other neurodegenerative diseases, with promising results in preclinical studies.

\textbf{Cancer:} Autophagy has a complex and context-dependent role in cancer. In early stages of tumor development, autophagy acts as a tumor suppressor by preventing the accumulation of damaged organelles and proteins that could promote genomic instability and malignant transformation. However, in established tumors, autophagy can promote tumor survival by providing nutrients and energy to rapidly growing cancer cells and by protecting them from the stresses of the tumor microenvironment, including hypoxia, nutrient deprivation, and chemotherapy. The modulation of autophagy---either activation or inhibition---is being explored as a therapeutic strategy for various types of cancer.

\textbf{Infectious disease:} Autophagy plays a critical role in innate immune defense against intracellular pathogens, including bacteria, viruses, and parasites. The process of ``xenophagy''---the autophagic degradation of intracellular microbes---has been shown to be an important mechanism for controlling infections. Defects in xenophagy have been linked to increased susceptibility to tuberculosis, Salmonella infections, and other intracellular infections. Autophagy also plays a role in the activation of adaptive immune responses by promoting antigen presentation on MHC molecules.

\textbf{Metabolic disease:} Autophagy is essential for the maintenance of metabolic homeostasis, particularly in tissues with high metabolic demands such as the liver, heart, and skeletal muscle. Defects in autophagy have been implicated in the pathogenesis of type 2 diabetes, obesity, and non-alcoholic fatty liver disease. The activation of autophagy through caloric restriction, exercise, or pharmacological agents has been shown to improve metabolic health in animal models and is being explored as a therapeutic strategy for metabolic diseases.

\textbf{Aging:} Autophagy declines with age, and this decline has been proposed to contribute to the accumulation of cellular damage that characterizes aging. The activation of autophagy through caloric restriction, rapamycin, or other interventions has been shown to extend lifespan in model organisms, including yeast, worms, flies, and mice. These findings have stimulated interest in the potential of autophagy activation as a strategy for promoting healthy aging and extending lifespan in humans.

\textbf{Therapeutic development:} The understanding of autophagy mechanisms has enabled the development of pharmacological modulators of autophagy. Rapamycin and its analogs (rapalogs) inhibit mTOR and thereby activate autophagy, and are being tested in clinical trials for various diseases. Other autophagy modulators, including trehalose, spermidine, and hydroxychloroquine (an autophagy inhibitor), are also being explored for therapeutic applications.

Yoshinori Ohsumi was awarded the Nobel Prize in Physiology or Medicine in 2016 ``for his discoveries of mechanisms for autophagy.'' His work has fundamentally changed our understanding of cellular physiology and has had a lasting impact on medicine.

\section{Legacy}

Yoshinori Ohsumi's legacy is embodied in the fundamental understanding of autophagy that his work established and the countless therapeutic avenues that have been opened by his discoveries. His identification of the first autophagy genes in yeast provided the foundation for a field that has grown to encompass thousands of researchers worldwide and has become one of the most active and productive areas of cell biology.

Ohsumi's approach exemplifies the power of simple model organisms to illuminate fundamental biological processes. By studying autophagy in yeast---a single-celled organism---Ohsumi was able to identify conserved mechanisms that apply to virtually all eukaryotic cells, including human cells. His work demonstrates the importance of basic research and the value of genetic approaches in biology.

The field of autophagy research continues to expand rapidly, with new discoveries being made about the mechanisms of autophagy, its regulation, and its roles in health and disease. The development of autophagy modulators as therapeutic agents for neurodegenerative diseases, cancer, metabolic diseases, and aging represents one of the most exciting areas of drug development today.

Ohsumi's work has also inspired a broader appreciation of the importance of cellular quality control mechanisms in maintaining health and preventing disease. The understanding that cells possess intrinsic mechanisms for degrading damaged components and recycling cellular materials has changed the way scientists think about cellular health and has opened up new approaches to the treatment of diseases characterized by the accumulation of damaged or misfolded proteins.

Today, autophagy research continues to yield new insights into cellular physiology and disease, and the therapeutic potential of autophagy modulation is being actively explored in clinical trials. The legacy of Yoshinori Ohsumi's work provides the foundation for these ongoing investigations and will continue to shape the future of cell biology and medicine for generations to come.


\section{Modern Developments}

Ohsumi's autophagy work has led to modern cell biology and drug development. Understanding of autophagy has informed development of mTOR inhibitors (everolimus, temsirolimus) for cancer and aging. Understanding of selective autophagy has led to research on neurodegenerative diseases (Alzheimer's, Parkinson's). Understanding of autophagy in cancer has revealed complex roles---both tumor suppression and tumor promotion. Autophagy modulators are in clinical trials for cancer, aging, and infectious diseases.
